\section{Ablations}
\label{sec:ablations}

We measure component contributions by removing stages from the frozen LoCoMo configuration with the entity graph, using the seven test conversations. The main metric is recall at $k$: the mean share of gold turns, which the benchmark marks as holding the answer, in the top $k$. The flat-fusion baseline adds a BM25 keyword score and an embedding score at fixed weights, with no later ranking stages. The two cross-encoders are local models that each read a question and a turn together and produce a relevance score.

Removing both cross-encoders produces the largest measured reduction in recall at 10: 6.1 [+4.5, +7.8] points (Figure~\ref{fig:ablation}). Removing contiguity costs 2.5 [+1.4, +3.7] points, and removing the speaker prior costs 0.9 [+0.1, +1.8]. The ranker with the entity graph exceeds the flat fusion by 17.5 [+15.0, +20.2] points at 10, 20.9 at 5 and 9.2 at 50 (Table~\ref{tab:ladder}, Figure~\ref{fig:recall}). Cross-encoders help most near the top of the ranking: removing them costs 8.4 points at 5 and 1.2 at 50.

The larger cross-encoder and candidate depth are not separately isolated in the comparison with the first-round chain. That chain reranks the top 50 turns with MiniLM alone. The later ranker gains 4.0 [+2.4, +5.3] points, but it changes both the model and the candidate depth. Within the later ensemble, removing MiniLM changes recall at 10 by $-$0.1 [$-$0.9, +0.6], so the smaller model could be dropped without measurable loss.

Reranking is also evaluated in the secondary retrievers. In the answer runs, 5r is 1.8 points above 5a with the same reader and judge: 96.0 [94.7, 97.4] against 94.2 [92.9, 95.7], with overlapping intervals. On \lmes{}, a MiniLM rerank of the top 12 sessions raises test \rall{3} to 87.3. Its pooled gain is +1.00 [$-$0.03, +2.13]; another seed and 4,000 draws give [+0.03, +2.14]. We therefore report this result as secondary and claim no gain.

The deployed store omits the entity graph because its measured contribution is small. On LoCoMo, the entity graph walk, the entity field and backward contiguity together add 0.9 to 1.9 points of turn recall. Removing personalized PageRank alone changes recall at 10 by $-$0.2 [$-$0.9, +0.5], which does not rule out a gain under one point. Removing PageRank, the entity field and backward contiguity together lowers recall by 1.2 [+0.4, +2.0] points. Backward contiguity is not an entity link, and the three components were not removed one at a time, so this combined difference cannot be attributed to the graph alone.

On \lmes{}, the three components add nothing we could measure. In development, PageRank at weight 0.1 changes the pooled objective by $-$0.05 [$-$1.06, +0.96], and adding entity, fact or date terms to the BM25 text changes it by 0 to $-$0.81. One possible explanation, which we did not test, is that the dense channel already matches the paraphrases an entity link would bridge.

\begin{table}[H]
\caption{Ablation ladder on LoCoMo test, text only, recall at 10 in deployment order. The third row is the graph-free ranker, whose three changes are not separated on test. The fourth row is the frozen configuration with the graph components. The last row has no interval because no paired run of it against the flat fusion was made.}
\label{tab:ladder}
\small
\centering
\setlength{\tabcolsep}{3pt}
\begin{tabular}{@{}L{44mm}cl@{}}
\toprule
Configuration & @10 & $\Delta$@10 vs.\ flat fusion \\
\midrule
Flat fusion & 64.0 & \\
+ forward contiguity & 68.5 & +4.5 [+1.6, +7.2] \\
+ BM25 weight 0.25, speaker prior, rerank & 80.3 & +16.3 [+13.8, +19.2] \\
+ PageRank, entity field, backward contiguity & 81.5 & +17.5 [+15.0, +20.2] \\
\midrule
Ranker without cross-encoders & 75.4 & +11.4 \\
\bottomrule
\end{tabular}
\end{table}

\begin{figure}[H]
\centering
\includegraphics[width=0.7\textwidth]{fig_ablation.pdf}
\caption{Change in LoCoMo turn recall at $k$ when a part is removed from the ranker with the entity graph (ablated minus full), text only, with 95\% cluster bootstrap intervals of the paired differences over the seven test conversations. Rows overlap: the contiguity row and the PageRank row both remove backward contiguity. In the MiniLM row, the weight of bge-reranker-v2-m3 rises to 0.5.}
\label{fig:ablation}
\end{figure}

\begin{figure}[H]
\centering
\includegraphics[width=\textwidth]{fig_recall.pdf}
\caption{LoCoMo retrieval. (a) The share of 1,536 questions whose gold turns are all among the turns handed to the answer model, one line per step added to the 5a and 5r retrievers. (b) The mean share of gold turns in the top $k$ on the seven test conversations, one line per ranker chain. Panels differ in metric and retriever.}
\label{fig:recall}
\end{figure}

The context-budget experiments measure a tradeoff between accuracy and reader input. Doubling BM25's cut from 20 to 40 turns raises answer accuracy by 3.5 points with GPT-6 Sol reading, from 82.3 to 85.8. Using an adaptive cut that starts at 10 turns, rather than always supplying 40, gives up 0.3 points of accuracy while using 40 percent of the tokens.

The retrieval experiments compare preprocessing, ranking and retrieval units. At a cut of 40 turns, cleaning and stemming the question raises every-gold-turn recall from 0.648 to 0.704; the dense embedder reaches 0.753, and the reranker reaches 0.789. Adding the neighbour on each side of every hit raises recall to 0.830 but sends about 96 turns, so we did not use that setting. Session-first retrieval lowers recall to 0.583. The bge-m3 embedder recalls less than bge-small on this data, 0.733 against 0.753, and takes longer to embed.

Two additional search changes were evaluated only on LoCoMo development. Pseudo-relevance feedback re-runs the search with words from the top hits. It gains 2.7 points on multi-hop recall, or two of 74 questions, below the registered bar of three. Date windows yield no gain. Neither change was evaluated on test.
